\chapter{Analyse critique des spécifications}
\label{chap:analyse-specs}

\section{Pourquoi ce chapitre existe}

Un ingénieur qui reçoit un cahier des charges le vérifie avant de coder. Cette vérification
n'est pas une critique de ses auteurs : c'est une étape normale et nécessaire du travail.

Les documents remis au début du stage totalisent environ \num{4200} lignes réparties en neuf
fichiers. Toutes leurs valeurs numériques dérivées ont été \emph{recalculées} indépendamment,
par un script écrit à cette seule fin. Cette démarche a mis en évidence
\textbf{onze anomalies} : erreurs de calcul, contradictions internes, données manquantes.

\begin{encadre}[title={Méthode adoptée face à une contradiction}]
Une règle de priorité entre sources a été fixée \emph{a priori}, avant d'examiner les cas,
pour éviter tout arbitrage de circonstance :

\begin{enumerate}
  \item le guide du procédé, qui se déclare lui-même référence définitive ;
  \item la référence du scénario réel, qui décrit le cas à reproduire ;
  \item le fichier de constantes ;
  \item la spécification du format de sortie ;
  \item les documents archivés, à valeur indicative.
\end{enumerate}

Et par-dessus tout : \textbf{la cohérence physique et dimensionnelle tranche}. Une source
qui viole la conservation de la masse ou l'homogénéité des unités a tort, quel que soit son
rang.
\end{encadre}

\section{Vue d'ensemble des anomalies}

\begin{table}[H]
\centering
\caption{Les onze anomalies identifiées}
\label{tab:anomalies}
\begin{tabular}{@{}clcl@{}}
\toprule
\textbf{N\textsuperscript{o}} & \textbf{Anomalie} & \textbf{Gravité} & \textbf{Type} \\
\midrule
A-01 & Production d'acide 54 mal calculée & Critique & calcul \\
A-02 & Bornes de stock d'acide 29 fausses & Critique & calcul \\
A-03 & Stock initial d'IR12 incohérent & Important & calcul \\
A-04 & Limite décanteur : entrée ou sortie ? & Important & contradiction \\
A-05 & Destination des boues de cocristallisation & Important & contradiction \\
A-06 & Hauteur minimale d'IR12 sous le volume mort & Mineur & incohérence \\
A-07 & Capacité des cuves d'acide 29 : trois valeurs & Important & contradiction \\
A-08 & Capacités d'échelon : t/h ou t/j ? & Mineur & unité \\
A-09 & Deux listes de lignes décadmiantes & Mineur & contradiction \\
A-10 & Fichier de profils qualité absent & Important & donnée manquante \\
A-11 & Format de sortie : 5 ou 7 feuilles ? & Cosmétique & contradiction \\
\bottomrule
\end{tabular}
\end{table}

Deux données essentielles sont par ailleurs purement et simplement absentes : la borne
maximale des transferts interzones, et le seuil quantifiant l'obligation de produire du
DEC\_CL.

\section{A-01 --- La production d'acide 54 est mal calculée}

\subsection{Constat}

La simulation détaillée fournie annonce une production totale de \SI{6636}{\tonne} par jour.
Le recalcul donne \SI{6981,7}{\tonne}, soit un écart de \SI{5,2}{\percent} sur le flux
principal du système.

\subsection{Preuve}

En appliquant la formule du document lui-même, $\Pi_m = \sum_e C_e h_e / 24$ :

\begin{table}[H]
\centering
\caption{Recalcul de la production d'acide 54, ligne par ligne (\tP/j)}
\label{tab:recalcul-p54}
\begin{tabular}{@{}lS[table-format=4.1]S[table-format=4.0]S[table-format=3.1]@{}}
\toprule
\textbf{Ligne} & {\textbf{Recalculé}} & {\textbf{Document}} & {\textbf{Écart}} \\
\midrule
14EXT & 1505.0 & 1365 & 140.0 \\
14AB  & 2134.2 & 2077 & 57.2 \\
14CD  & 1040.0 & 1040 & 0.0 \\
14XY  & 820.8  & 779  & 41.8 \\
14ZU  & 1481.7 & 1375 & 106.7 \\
\midrule
\textbf{Total} & \bfseries 6981.7 & \bfseries 6636 & \bfseries 345.7 \\
\bottomrule
\end{tabular}
\end{table}

Détaillons la ligne 14EXT :
\[
\frac{420 \times 24 + 420 \times 24 + 420 \times 24 + 420 \times 14}{24}
= \frac{\num{36120}}{24} = \num{1505}
\]
Le document annonce \num{1365}. Il manque exactement $420 \times 8/24 = 140$.

\begin{encadre}[title={Un indice révélateur}]
La ligne 14CD est la \emph{seule} correcte --- et c'est aussi la seule dont tous les
échelons marchent \SI{24}{\hour}. Toutes les erreurs concernent des lignes ayant au moins un
échelon à horaire réduit. Le document a donc mal appliqué le prorata horaire.
\end{encadre}

\subsection{Arbitrage}

Recalculer systématiquement, ne recopier aucun chiffre. Cette règle est appliquée dans tout
le code : aucune valeur dérivée n'y est écrite en dur, et chaque formule est couverte par un
test unitaire.

\section{A-02 --- Les bornes de stock d'acide 29 sont surestimées de 59\,\%}

\subsection{Constat}

Le fichier de constantes annonce $Z^{\min}_{29} \approx \num{364,5}$ et
$Z^{\max}_{29} \approx \num{1530,7}$. Le recalcul de \emph{sa propre formule} donne
\num{229,8} et \num{965,3}.

\subsection{Validation croisée --- la preuve décisive}

La formule fournie est $Z^{\max} = 2 \times S \times h_{\max} \times \theta \times \rho$.
Appliquée à l'acide 54 :
\[
2 \times 95 \times 9{,}40 \times 0{,}50 \times 1{,}66 = \num{1482,4}
\]
à comparer aux \num{1483,6} annoncés par le document : \textbf{concordance au dixième
près}.

\begin{resultat}[title={Conclusion}]
La méthode de calcul est correcte --- elle reproduit exactement les valeurs du document pour
l'acide 54. L'erreur est donc \emph{localisée} sur l'acide 29.
\end{resultat}

\subsection{Signature de l'erreur}

Les deux valeurs erronées présentent le même facteur :
\[
\frac{364{,}5}{229{,}8} = 1{,}5860
\qquad
\frac{1530{,}7}{965{,}3} = 1{,}5858
\]
Un facteur constant sur les deux bornes indique une erreur de saisie unique --- un
coefficient mal recopié --- et non deux erreurs indépendantes.

\subsection{Arbitrage et impact}

Les valeurs recalculées sont retenues. Un argument supplémentaire l'impose : les stocks
\emph{initiaux} sont produits par la fonction de conversion. Comparer un stock calculé par
cette fonction à une borne obtenue autrement serait dimensionnellement incohérent.

L'impact est majeur. Avec la borne correcte (\num{965} au lieu de \num{1531}), plusieurs
lignes atteignent leur capacité et les transferts interzones deviennent
\textbf{indispensables} --- ce qui donne rétrospectivement tout son sens à l'existence de la
matrice de transferts dans le cahier des charges.

\section{A-04 --- La limite des décanteurs porte-t-elle sur l'entrée ou la sortie ?}

Le guide du procédé se contredit. En un endroit, il précise
« \emph{Total Clarification Limit: 1000 t/day per line (before yield)} », donc en entrée.
Ailleurs, un exemple calcule « envoyer $1000/0{,}9 = 1111$ t en clarification $\to$ environ
\SI{1000}{\tonne} de DEC CL », ce qui place la limite en sortie. Les règles de validation du
même dossier vérifient quant à elles une quantité \emph{envoyée}.

\paragraph{Arbitrage.} La limite porte sur l'entrée, pour trois raisons convergentes :
c'est la lecture physique --- un décanteur est dimensionné par le débit qu'il reçoit, il
s'agit d'un temps de séjour ; la mention « \emph{before yield} » est explicite ; et deux
sources sur trois concordent.

\section{A-05 --- Où retournent les boues de cocristallisation ?}

Six sources, deux réponses opposées.

\begin{table}[H]
\centering
\caption{Sources contradictoires sur la destination des boues de cocristallisation}
\label{tab:boues-coc}
\begin{tabular}{@{}ll@{}}
\toprule
\textbf{Source} & \textbf{Destination indiquée} \\
\midrule
Guide du procédé (référence définitive) & stock d'acide \textbf{29} standard \\
Référence du scénario réel & stock d'acide \textbf{29} standard \\
Simulation détaillée & stock d'acide \textbf{29} standard \\
Fichier de constantes & stock d'acide \textbf{54} NCL \\
Spécification du format de sortie & stock d'acide \textbf{54} NCL \\
Document archivé & stock d'acide \textbf{54} NCL \\
\bottomrule
\end{tabular}
\end{table}

\paragraph{Arbitrage.} Retour au stock d'acide 29 standard, conformément à la hiérarchie des
sources : le guide se déclare référence définitive, la référence du scénario le confirme, et
la simulation applique effectivement ce choix. Deux des trois sources contraires sont
archivées.

\begin{encadre}[title={Un argument physique vient confirmer l'arbitrage documentaire}]
Les eaux mères de cristallisation sont chargées d'impuretés. Les renvoyer au niveau 54
réintroduirait ces impuretés dans un acide déjà concentré. Les diluer dans la masse de
l'acide 29, en amont de toute la chaîne de purification, est plus sain.

\emph{Le choix physiquement le plus plausible coïncide avec le choix documentairement le
mieux appuyé.} Le point de retour est néanmoins resté paramétrable dans le code, pour
permettre une bascule immédiate si l'encadrant tranche autrement.
\end{encadre}

\section{A-08 --- Les capacités d'échelon sont-elles en t/h ou en t/j ?}

Le fichier de constantes annote les capacités « tonnes/hour », tout en fournissant la
formule $\Pi = C \times h / 24$.

\begin{proof}[Preuve par l'absurde]
Si $C$ était en \si{\tonne\per\hour}, un échelon de \SI{420}{\tonne\per\hour} marchant
\SI{24}{\hour} produirait \SI{10080}{\tonne} par jour, et la division par 24 n'aurait aucun
sens dimensionnel. Cinq lignes à ce régime produiraient plus de
\SI{50000}{\tonne}~\PdeuxOcinq{} par jour, soit près de vingt fois la production quotidienne
du groupe.
\end{proof}

Les capacités sont donc en tonnes par jour à plein régime, et la formule est un simple
prorata horaire. L'impact sur les calculs est nul --- la formule était déjà appliquée
correctement --- mais la clarification importe pour la documentation et le code.

\section{A-10 --- Le fichier des profils qualité est absent}

Ce fichier est qualifié d'essentiel dans quatre documents, mais ne figure pas dans le
dossier remis. Ses coefficients ont dû être reconstitués un à un à partir du texte des
documents.

\begin{table}[H]
\centering
\caption{Profils qualité reconstruits (tonnes d'acide par tonne d'engrais)}
\label{tab:profils}
\begin{tabular}{@{}llS[table-format=1.3]@{}}
\toprule
\textbf{Produit} & \textbf{Type d'acide} & {\textbf{Coefficient}} \\
\midrule
DAP\_STANDARD & acide 29 std & 0.127 \\
DAP\_STANDARD & acide 54 NCL & 0.354 \\
TSP\_EURO & acide 54 dec total & 0.379 \\
MAP\_11\_52\_EU & acide 29 dec & 0.030 \\
MAP\_11\_52\_EU & acide 54 dec total & 0.129 \\
NPK\_12\_24\_12\_EU & acide 54 dec total & 0.091 \\
NPK\_15\_15\_15\_EU & acide 29 dec & 0.171 \\
\bottomrule
\end{tabular}
\end{table}

Le fichier reconstruit indique explicitement, pour chaque coefficient, le document dont il
provient, et porte la mention « reconstruit » afin d'être remplacé sans ambiguïté dès que
l'original sera disponible.

\begin{encadre}[title={Alerte supplémentaire}]
Un même produit apparaît sous \emph{deux noms différents} dans le même scénario :
\texttt{MAP\_11\_54\_NE\_EU} et \texttt{MAP\_11\_52\_EU}. Il s'agit vraisemblablement d'une
coquille. Le second est retenu, cohérent avec la nomenclature commerciale MAP 11-52-0.
\end{encadre}

\section{Les deux données manquantes}

\subsection{La borne maximale des transferts}

Le script de test fourni importe et utilise une constante \texttt{MAX\_TRANSFER}, mais le
fichier de constantes ne définit que les minima. La donnée existe donc dans le code
d'origine sans avoir été transcrite dans la documentation.

\subsection{Le seuil de production obligatoire de DEC\_CL}

C'est le seul endroit du dossier où une règle est qualifiée d'\emph{obligatoire} sans être
chiffrée. Or elle est décisive : sans seuil, l'optimiseur ne produira jamais de DEC\_CL et
violera la règle métier --- exactement le risque que le document anticipe lui-même.

\begin{encadre}[title={La démarche adoptée face à une donnée manquante}]
Ni deviner en silence, ni s'arrêter, mais :
\begin{enumerate}
  \item poser une valeur d'attente \textbf{explicitement signalée comme telle} ;
  \item la rendre \textbf{paramétrable}, jamais écrite en dur ;
  \item \textbf{poser la question} à qui détient la réponse ;
  \item \textbf{mesurer sa sensibilité} par une analyse dédiée.
\end{enumerate}

Le chapitre~\ref{chap:resultats} montre que cette démarche a permis, pour l'une des deux
données, de \emph{résoudre la question sans jamais obtenir la réponse}.
\end{encadre}

\section{Enseignements}

\subsection{Ne jamais recopier un chiffre --- le recalculer}

Quatre des onze anomalies sont de simples erreurs arithmétiques. Toutes ont été détectées en
réappliquant les formules du document à ses propres données.

\emph{Conséquence pour le code} : aucune valeur dérivée n'y est écrite en dur. Tout ce qui
peut être calculé l'est, à partir des données primaires, et un test unitaire vérifie chaque
formule.

\subsection{Verrouiller un constat par un test}

Un constat documenté mais non testé finit par être oublié. Chaque anomalie arbitrée a donc
donné lieu à un test automatisé qui échouera si la valeur retenue change :

\begin{lstlisting}[language=Python, caption={Test verrouillant l'anomalie A-02}, label={lst:test-a02}]
def test_bornes_29_different_du_dossier_anomalie_a02():
    """Les bornes recalculees s'ecartent du dossier d'un facteur 1,586."""
    z_min, z_max = units.bornes_stock_29()
    assert z_max == pytest.approx(965.3, abs=0.5)
    assert 1530.7 / z_max == pytest.approx(1.586, abs=0.01)
\end{lstlisting}

Si l'encadrant fournit un jour d'autres valeurs, ce test échouera et signalera qu'il faut
réexaminer l'arbitrage.
